% =============================================================
%  02_esercizi.tex — Tutti gli esercizi (148 totali, 10 capitoli)
%  \part{Esercizi} e la legenda sono dichiarati in main.tex
% =============================================================

% =============================================================
\setcounter{section}{0}
\section{Fondamenti}
% =============================================================

\noindent
In questo capitolo si prende confidenza con la sintassi di
Common Lisp: S-expression, \fn{defun}, aritmetica, predicati,
condizionali (\fn{if}, \fn{cond}), binding locale (\fn{let},
\fn{let*}) e iterazione base (\fn{dotimes}, \fn{loop}).

\bigskip

\label{ex:1}
\begin{essbox}{1}{Somma}{1}
\exlabel{Implementare:}
\nopagebreak
\begin{lstlisting}
(add a b)
\end{lstlisting}
\begin{lstlisting}
(add 3 5)      ; -> 8
(add -2 7)     ; -> 5
(add 1.5 2.5)  ; -> 4.0
\end{lstlisting}


\small\color{black!55}\hyperref[sol:1]{$\rightarrow$ Soluzione}\normalsize\color{black}
\end{essbox}

\label{ex:2}
\begin{stdbox}{2}{Geometria del cerchio}{1}
\exlabel{Implementare:}
\nopagebreak
\begin{lstlisting}
(circle-area radius)
(circle-circumference radius)
\end{lstlisting}
Usare la costante \fn{pi} per $\pi$.
\begin{lstlisting}
(circle-area 2)           ; -> 12.566...
(circle-circumference 2)  ; -> 12.566...
\end{lstlisting}


\small\color{black!55}\hyperref[sol:2]{$\rightarrow$ Soluzione}\normalsize\color{black}
\end{stdbox}

\label{ex:3}
\begin{stdbox}{3}{Conversione Celsius/Fahrenheit}{1}
\exlabel{Implementare:}
\nopagebreak
\begin{lstlisting}
(celsius->fahrenheit c)
(fahrenheit->celsius f)
\end{lstlisting}
Formule: $F = C \times 9/5 + 32$;\quad $C = (F - 32) \times 5/9$.


\small\color{black!55}\hyperref[sol:3]{$\rightarrow$ Soluzione}\normalsize\color{black}
\end{stdbox}

\label{ex:4}
\begin{essbox}{4}{Ipotenusa}{1}
\exlabel{Implementare:}
\nopagebreak
\begin{lstlisting}
(ipotenusa a b)
\end{lstlisting}
\begin{lstlisting}
(ipotenusa 3 4)   ; -> 5.0
\end{lstlisting}
\vincolo{usare \fn{let} per i quadrati intermedi.}


\small\color{black!55}\hyperref[sol:4]{$\rightarrow$ Soluzione}\normalsize\color{black}
\end{essbox}

\label{ex:5}
\begin{colbox}{5}{Distanza euclidea}{1}
\exlabel{Implementare:}
\nopagebreak
\begin{lstlisting}
(distanza-euclidea x1 y1 x2 y2)
\end{lstlisting}
\begin{lstlisting}
(distanza-euclidea 0 0 3 4)   ; -> 5.0
\end{lstlisting}
\vincolo{usare \fn{let} per i valori intermedi.}


\small\color{black!55}\hyperref[sol:5]{$\rightarrow$ Soluzione}\normalsize\color{black}
\end{colbox}

\label{ex:6}
\begin{essbox}{6}{Pari o dispari}{1}
\exlabel{Implementare:}
\nopagebreak
\begin{lstlisting}
(pari-p n)
\end{lstlisting}
senza usare \fn{evenp}.
\begin{lstlisting}
(pari-p 8)    ; -> T
(pari-p 7)    ; -> NIL
\end{lstlisting}


\small\color{black!55}\hyperref[sol:6]{$\rightarrow$ Soluzione}\normalsize\color{black}
\end{essbox}

\label{ex:7}
\begin{colbox}{7}{Divisibilita'}{1}
\exlabel{Implementare:}
\nopagebreak
\begin{lstlisting}
(divisibile-p a b)
\end{lstlisting}
\begin{lstlisting}
(divisibile-p 12 4)   ; -> T
(divisibile-p 7  3)   ; -> NIL
\end{lstlisting}


\small\color{black!55}\hyperref[sol:7]{$\rightarrow$ Soluzione}\normalsize\color{black}
\end{colbox}

\label{ex:8}
\begin{essbox}{8}{Segno}{1}
\exlabel{Implementare:}
\nopagebreak
\begin{lstlisting}
(segno n)
\end{lstlisting}
che restituisce \fn{'positivo}, \fn{'negativo} o \fn{'zero}.
\begin{lstlisting}
(segno -7)   ; -> NEGATIVO
(segno 0)    ; -> ZERO
(segno 12)   ; -> POSITIVO
\end{lstlisting}
\vincolo{usare \fn{cond}.}


\small\color{black!55}\hyperref[sol:8]{$\rightarrow$ Soluzione}\normalsize\color{black}
\end{essbox}

\label{ex:9}
\begin{stdbox}{9}{Massimo di tre}{1}
\exlabel{Implementare} \fn{(max-tre a b c)} senza usare \fn{max}.
\nopagebreak
\begin{lstlisting}
(max-tre 3 7 2)   ; -> 7
\end{lstlisting}
\vincolo{usare \fn{cond}.}


\small\color{black!55}\hyperref[sol:9]{$\rightarrow$ Soluzione}\normalsize\color{black}
\end{stdbox}

\label{ex:10}
\begin{colbox}{10}{Triangolo valido}{1}
\exlabel{Implementare} \fn{(triangolo-valido-p a b c)}.
Ogni lato deve essere minore della somma degli altri due.
\begin{lstlisting}
(triangolo-valido-p 3 4 5)   ; -> T
(triangolo-valido-p 1 2 10)  ; -> NIL
\end{lstlisting}


\small\color{black!55}\hyperref[sol:10]{$\rightarrow$ Soluzione}\normalsize\color{black}
\end{colbox}

\label{ex:11}
\begin{essbox}{11}{Anno bisestile}{2}
\exlabel{Implementare:}
\nopagebreak
\begin{lstlisting}
(bisestile-p anno)
\end{lstlisting}
Divisibile per 4, eccetto multipli di 100 che non siano anche di 400.
\begin{lstlisting}
(bisestile-p 2000)   ; -> T
(bisestile-p 1900)   ; -> NIL
(bisestile-p 2024)   ; -> T
\end{lstlisting}
\vincolo{usare \fn{and}/\fn{or} senza \fn{cond}.}


\small\color{black!55}\hyperref[sol:11]{$\rightarrow$ Soluzione}\normalsize\color{black}
\end{essbox}

\label{ex:12}
\begin{stdbox}{12}{Giorni del mese}{2}
\exlabel{Implementare} \fn{(giorni-mese mese anno)} con \fn{mese} da 1 a 12.
\nopagebreak
\begin{lstlisting}
(giorni-mese 2 2024)   ; -> 29
(giorni-mese 2 2023)   ; -> 28
(giorni-mese 7 2024)   ; -> 31
\end{lstlisting}
Usare \fn{bisestile-p} per febbraio.


\small\color{black!55}\hyperref[sol:12]{$\rightarrow$ Soluzione}\normalsize\color{black}
\end{stdbox}

\label{ex:13}
\begin{colbox}{13}{Voto in lettere}{1}
\exlabel{Implementare} \fn{(voto->giudizio punteggio)}, punteggio 0--100.

90--100 $\to$ ``ottimo''; 75--89 $\to$ ``buono'';
60--74 $\to$ ``sufficiente''; $<60$ $\to$ ``insufficiente''.

\vincolo{usare \fn{cond}.}


\small\color{black!55}\hyperref[sol:13]{$\rightarrow$ Soluzione}\normalsize\color{black}
\end{colbox}

\label{ex:14}
\begin{colbox}{14}{Terna pitagorica}{1}
\exlabel{Implementare} \fn{(terna-pitagorica-p a b c)}.
\nopagebreak
\begin{lstlisting}
(terna-pitagorica-p 3 4 5)    ; -> T
(terna-pitagorica-p 1 2 3)    ; -> NIL
\end{lstlisting}


\small\color{black!55}\hyperref[sol:14]{$\rightarrow$ Soluzione}\normalsize\color{black}
\end{colbox}

\label{ex:15}
\begin{essbox}{15}{Radici di equazione di secondo grado}{2}
\exlabel{Implementare:}
\nopagebreak
\begin{lstlisting}
(radici-quadratiche a b c)
\end{lstlisting}
Restituire lista con le due radici reali, oppure \fn{NIL} se $\Delta < 0$.
\begin{lstlisting}
(radici-quadratiche 1 -5 6)   ; -> (3.0 2.0)
(radici-quadratiche 1  0 1)   ; -> NIL
\end{lstlisting}
\vincolo{usare \fn{let*} per discriminante e radici.}


\small\color{black!55}\hyperref[sol:15]{$\rightarrow$ Soluzione}\normalsize\color{black}
\end{essbox}

\label{ex:16}
\begin{stdbox}{16}{Calcolatrice simbolica}{2}
\exlabel{Implementare} \fn{(calcola op a b)} dove \fn{op} e' un simbolo
\fn{'+}, \fn{'-}, \fn{'*}, \fn{'/}.
\begin{lstlisting}
(calcola '+ 3 4)   ; -> 7
(calcola '* 5 6)   ; -> 30
\end{lstlisting}
Gestire la divisione per zero restituendo \fn{NIL}.


\small\color{black!55}\hyperref[sol:16]{$\rightarrow$ Soluzione}\normalsize\color{black}
\end{stdbox}

\label{ex:17}
\begin{essbox}{17}{Somma delle cifre}{2}
\exlabel{Implementare:}
\nopagebreak
\begin{lstlisting}
(somma-cifre n)
\end{lstlisting}
\begin{lstlisting}
(somma-cifre 1234)   ; -> 10
(somma-cifre 99)     ; -> 18
\end{lstlisting}
\textbf{Suggerimento:} estrarre le cifre con \fn{mod} e \fn{floor}.


\small\color{black!55}\hyperref[sol:17]{$\rightarrow$ Soluzione}\enspace\textcolor{black!30}{|}\enspace\hyperref[hint:17]{$\rightarrow$ Suggerimento}\normalsize\color{black}
\end{essbox}

\label{ex:18}
\begin{essbox}{18}{Conta fino a N}{1}
\exlabel{Implementare} \fn{(conta-fino-a n)} che stampa i numeri da 1 a \fn{n}.
\nopagebreak
\begin{lstlisting}
(conta-fino-a 4)
; 1
; 2
; 3
; 4
\end{lstlisting}
\vincolo{usare \fn{dotimes} o \fn{loop}.}


\small\color{black!55}\hyperref[sol:18]{$\rightarrow$ Soluzione}\normalsize\color{black}
\end{essbox}

\label{ex:19}
\begin{colbox}{19}{Tabellina}{1}
\exlabel{Implementare} \fn{(tabellina n)} che stampa la tabellina di \fn{n}
da 1 a 10.
\begin{lstlisting}
(tabellina 3)
; 3 x 1 = 3
; 3 x 2 = 6
; ...
\end{lstlisting}


\small\color{black!55}\hyperref[sol:19]{$\rightarrow$ Soluzione}\normalsize\color{black}
\end{colbox}

\label{ex:20}
\begin{colbox}{20}{Sconto}{1}
\exlabel{Implementare} \fn{(sconto prezzo)}:
$\geq 200$ $\to$ 20\%; $\geq 100$ $\to$ 10\%; altrimenti 0\%.
\begin{lstlisting}
(sconto 250)   ; -> 200.0
(sconto 120)   ; -> 108.0
(sconto 50)    ; -> 50
\end{lstlisting}


\small\color{black!55}\hyperref[sol:20]{$\rightarrow$ Soluzione}\normalsize\color{black}
\end{colbox}

\label{ex:21}
\begin{colbox}{21}{Paga settimanale}{1}
\exlabel{Implementare} \fn{(paga-settimanale ore tariffa)}.
Fino a 40 ore: tariffa normale.
Ore eccedenti: tariffa $\times$ 1.5.
\begin{lstlisting}
(paga-settimanale 35 10)   ; -> 350
(paga-settimanale 45 10)   ; -> 475
\end{lstlisting}


\small\color{black!55}\hyperref[sol:21]{$\rightarrow$ Soluzione}\normalsize\color{black}
\end{colbox}

% =============================================================
\setcounter{section}{1}
\section{Liste}
% =============================================================

\noindent
Le liste sono la struttura dati fondamentale di Common Lisp.
In questo capitolo si implementano da zero le operazioni
principali usando \fn{car}, \fn{cdr}, \fn{cons} e la ricorsione.
\textbf{Non usare} le funzioni built-in che risolvono direttamente
l'esercizio, salvo diversa indicazione.

\bigskip

\label{ex:22}
\begin{essbox}{22}{Costruire liste con cons e list}{1}
\textbf{Parte A.} Costruire \fn{'(a (b c) d)} usando solo \fn{cons}.

\textbf{Parte B.} Costruire la stessa lista con \fn{list}.

Verificare che le due espressioni siano \fn{equal}.


\small\color{black!55}\hyperref[sol:22]{$\rightarrow$ Soluzione}\normalsize\color{black}
\end{essbox}

\label{ex:23}
\begin{stdbox}{23}{Accesso con car e cdr}{1}
Data \fn{'(a b c d e)}, estrarre il terzo elemento
usando solo \fn{car} e \fn{cdr} (non \fn{nth}, \fn{third}, ecc.).
\begin{lstlisting}
; risultato atteso: C
\end{lstlisting}


\small\color{black!55}\hyperref[sol:23]{$\rightarrow$ Soluzione}\normalsize\color{black}
\end{stdbox}

\label{ex:24}
\begin{essbox}{24}{Primo e resto}{1}
\exlabel{Implementare:}
\nopagebreak
\begin{lstlisting}
(first-element lst)
(rest-elements lst)
\end{lstlisting}
Gestire la lista vuota restituendo \fn{NIL}.
\begin{lstlisting}
(first-element '(a b c))    ; -> A
(rest-elements '(a b c))    ; -> (B C)
(first-element '())         ; -> NIL
\end{lstlisting}


\small\color{black!55}\hyperref[sol:24]{$\rightarrow$ Soluzione}\normalsize\color{black}
\end{essbox}

\label{ex:25}
\begin{colbox}{25}{Primo e ultimo}{1}
\exlabel{Implementare} \fn{(primo-e-ultimo lst)} che restituisce una lista
con il primo e l'ultimo elemento.
\begin{lstlisting}
(primo-e-ultimo '(a b c d))   ; -> (A D)
(primo-e-ultimo '(a))         ; -> (A A)
\end{lstlisting}
Gestire la lista vuota.


\small\color{black!55}\hyperref[sol:25]{$\rightarrow$ Soluzione}\normalsize\color{black}
\end{colbox}

\label{ex:26}
\begin{colbox}{26}{Scambia i primi due}{1}
\exlabel{Implementare} \fn{(scambia-primi-due lst)}.
\nopagebreak
\begin{lstlisting}
(scambia-primi-due '(a b c d))   ; -> (B A C D)
\end{lstlisting}
Gestire liste con meno di due elementi.


\small\color{black!55}\hyperref[sol:26]{$\rightarrow$ Soluzione}\normalsize\color{black}
\end{colbox}

\label{ex:27}
\begin{essbox}{27}{Lunghezza}{2}
\exlabel{Implementare} \fn{(length* lst)} senza usare \fn{length}.
\nopagebreak
\begin{lstlisting}
(length* '())        ; -> 0
(length* '(a b c))   ; -> 3
\end{lstlisting}
\vincolo{ricorsione.}


\small\color{black!55}\hyperref[sol:27]{$\rightarrow$ Soluzione}\normalsize\color{black}
\end{essbox}

\label{ex:28}
\begin{stdbox}{28}{Somma degli elementi}{2}
\exlabel{Implementare} \fn{(sum-list lst)}.
\nopagebreak
\begin{lstlisting}
(sum-list '(1 2 3 4))   ; -> 10
(sum-list '())          ; -> 0
\end{lstlisting}
\vincolo{ricorsione.}


\small\color{black!55}\hyperref[sol:28]{$\rightarrow$ Soluzione}\normalsize\color{black}
\end{stdbox}

\label{ex:29}
\begin{colbox}{29}{Prodotto degli elementi}{2}
\exlabel{Implementare} \fn{(product-list lst)}.
\nopagebreak
\begin{lstlisting}
(product-list '(2 3 4))   ; -> 24
(product-list '())        ; -> 1
\end{lstlisting}


\small\color{black!55}\hyperref[sol:29]{$\rightarrow$ Soluzione}\normalsize\color{black}
\end{colbox}

\label{ex:30}
\begin{colbox}{30}{Massimo di una lista}{2}
\exlabel{Implementare} \fn{(max-list lst)}.
\nopagebreak
\begin{lstlisting}
(max-list '(4 8 2 9 1))   ; -> 9
\end{lstlisting}
Definire il comportamento per la lista vuota.


\small\color{black!55}\hyperref[sol:30]{$\rightarrow$ Soluzione}\normalsize\color{black}
\end{colbox}

\label{ex:31}
\begin{essbox}{31}{Appartenenza}{2}
\exlabel{Implementare} \fn{(member* item lst)} senza usare \fn{member}.
\nopagebreak
\begin{lstlisting}
(member* 'b '(a b c))   ; -> T
(member* 'x '(a b c))   ; -> NIL
\end{lstlisting}


\small\color{black!55}\hyperref[sol:31]{$\rightarrow$ Soluzione}\normalsize\color{black}
\end{essbox}

\label{ex:32}
\begin{stdbox}{32}{Conta occorrenze}{2}
\exlabel{Implementare} \fn{(count-element item lst)} senza usare \fn{count}.
\nopagebreak
\begin{lstlisting}
(count-element 'a '(a b a c a))   ; -> 3
\end{lstlisting}


\small\color{black!55}\hyperref[sol:32]{$\rightarrow$ Soluzione}\normalsize\color{black}
\end{stdbox}

\label{ex:33}
\begin{colbox}{33}{Posizione}{2}
\exlabel{Implementare} \fn{(position* item lst)}: indice 0-based
della prima occorrenza, oppure \fn{NIL}.
\begin{lstlisting}
(position* 'c '(a b c d))   ; -> 2
(position* 'x '(a b c d))   ; -> NIL
\end{lstlisting}


\small\color{black!55}\hyperref[sol:33]{$\rightarrow$ Soluzione}\normalsize\color{black}
\end{colbox}

\label{ex:34}
\begin{colbox}{34}{Ultimo elemento}{2}
\exlabel{Implementare} \fn{(last-element lst)} senza usare \fn{last}.
\nopagebreak
\begin{lstlisting}
(last-element '(a b c))   ; -> C
\end{lstlisting}


\small\color{black!55}\hyperref[sol:34]{$\rightarrow$ Soluzione}\normalsize\color{black}
\end{colbox}

\label{ex:35}
\begin{essbox}{35}{Incrementa tutti}{2}
\exlabel{Implementare} \fn{(incrementa-tutti lst)} che aggiunge 1
a ogni elemento.
\begin{lstlisting}
(incrementa-tutti '(1 2 3))   ; -> (2 3 4)
\end{lstlisting}
\vincolo{ricorsione; non usare \fn{mapcar}.}

\exnota{Obiettivo: capire il pattern di trasformazione, precursore di \fn{mapcar}.}


\small\color{black!55}\hyperref[sol:35]{$\rightarrow$ Soluzione}\normalsize\color{black}
\end{essbox}

\label{ex:36}
\begin{stdbox}{36}{Filtra i pari}{2}
\exlabel{Implementare} \fn{(solo-pari lst)}.
\nopagebreak
\begin{lstlisting}
(solo-pari '(1 2 3 4 5 6))   ; -> (2 4 6)
\end{lstlisting}
\vincolo{ricorsione; non usare \fn{remove-if}.}

\exnota{Obiettivo: capire il pattern di filtraggio, precursore di \fn{remove-if-not}.}


\small\color{black!55}\hyperref[sol:36]{$\rightarrow$ Soluzione}\normalsize\color{black}
\end{stdbox}

\label{ex:37}
\begin{stdbox}{37}{Elimina tutte le occorrenze}{2}
\exlabel{Implementare} \fn{(elimina-tutte item lst)} senza usare \fn{remove}.
\nopagebreak
\begin{lstlisting}
(elimina-tutte 'a '(a b a c a))   ; -> (B C)
\end{lstlisting}


\small\color{black!55}\hyperref[sol:37]{$\rightarrow$ Soluzione}\normalsize\color{black}
\end{stdbox}

\label{ex:38}
\begin{colbox}{38}{Seleziona posizioni dispari}{2}
\exlabel{Implementare} \fn{(seleziona-dispari lst)}: restituisce gli elementi
in posizione 1\textsuperscript{a}, 3\textsuperscript{a}, 5\textsuperscript{a},
\ldots\ nella numerazione naturale (indici 0, 2, 4, \ldots\ in notazione 0-based).
\begin{lstlisting}
(seleziona-dispari '(a b c d e))   ; -> (A C E)
\end{lstlisting}


\small\color{black!55}\hyperref[sol:38]{$\rightarrow$ Soluzione}\normalsize\color{black}
\end{colbox}

\label{ex:39}
\begin{essbox}{39}{Reverse con accumulatore}{2}
\exlabel{Implementare} \fn{(reverse* lst)} senza usare \fn{reverse}.

Prima scrivere la versione ricorsiva semplice, poi quella
con accumulatore (tail-recursive).
\begin{lstlisting}
(reverse* '(a b c d))   ; -> (D C B A)
\end{lstlisting}


\small\color{black!55}\hyperref[sol:39]{$\rightarrow$ Soluzione}\normalsize\color{black}
\end{essbox}

\label{ex:40}
\begin{essbox}{40}{Append}{3}
\exlabel{Implementare} \fn{(append* lst1 lst2)} senza usare \fn{append}.
\nopagebreak
\begin{lstlisting}
(append* '(a b) '(c d))   ; -> (A B C D)
(append* '() '(a b))      ; -> (A B)
\end{lstlisting}
Non modificare le liste originali; ricorsione su \fn{lst1}.



\small\color{black!55}\hyperref[sol:40]{$\rightarrow$ Soluzione}\enspace\textcolor{black!30}{|}\enspace\hyperref[hint:40]{$\rightarrow$ Suggerimento}\normalsize\color{black}
\end{essbox}

\label{ex:41}
\begin{stdbox}{41}{Rimozione prima occorrenza}{3}
\exlabel{Implementare} \fn{(remove-first item lst)}: rimuove solo la
prima occorrenza.
\begin{lstlisting}
(remove-first 'b '(a b c b d))   ; -> (A C B D)
\end{lstlisting}
Non usare \fn{remove}; non modificare la lista originale.



\small\color{black!55}\hyperref[sol:41]{$\rightarrow$ Soluzione}\enspace\textcolor{black!30}{|}\enspace\hyperref[hint:41]{$\rightarrow$ Suggerimento}\normalsize\color{black}
\end{stdbox}

\label{ex:42}
\begin{colbox}{42}{Sostituzione prima occorrenza}{2}
\exlabel{Implementare} \fn{(replace-first old new lst)}.
\nopagebreak
\begin{lstlisting}
(replace-first 'b 'x '(a b c b))   ; -> (A X C B)
\end{lstlisting}


\small\color{black!55}\hyperref[sol:42]{$\rightarrow$ Soluzione}\normalsize\color{black}
\end{colbox}

\label{ex:43}
\begin{colbox}{43}{Lista palindroma}{2}
\exlabel{Implementare} \fn{(palindroma-p lst)}.
\nopagebreak
\begin{lstlisting}
(palindroma-p '(a b c b a))   ; -> T
(palindroma-p '(a b c))       ; -> NIL
\end{lstlisting}
Usare la funzione \fn{reverse*} dell'es.~39 (o \fn{reverse}) e \fn{equal}.


\small\color{black!55}\hyperref[sol:43]{$\rightarrow$ Soluzione}\normalsize\color{black}
\end{colbox}

\label{ex:44}
\begin{stdbox}{44}{Range}{2}
\exlabel{Implementare} \fn{(range a b)}: interi da \fn{a} a \fn{b} inclusi.
\nopagebreak
\begin{lstlisting}
(range 1 5)    ; -> (1 2 3 4 5)
(range 3 3)    ; -> (3)
(range 5 3)    ; -> NIL
\end{lstlisting}


\small\color{black!55}\hyperref[sol:44]{$\rightarrow$ Soluzione}\normalsize\color{black}
\end{stdbox}

\label{ex:45}
\begin{colbox}{45}{Duplica elementi}{2}
\exlabel{Implementare} \fn{(duplica-elementi lst)}.
\nopagebreak
\begin{lstlisting}
(duplica-elementi '(a b c))   ; -> (A A B B C C)
\end{lstlisting}


\small\color{black!55}\hyperref[sol:45]{$\rightarrow$ Soluzione}\normalsize\color{black}
\end{colbox}

\label{ex:46}
\begin{colbox}{46}{Estrai chiavi da alist}{2}
\exlabel{Implementare} \fn{(estrai-chiavi alist)}: restituisce la
lista delle chiavi (\fn{car} di ogni coppia puntata).
\begin{lstlisting}
(estrai-chiavi '((a . 1) (b . 2) (c . 3)))
; -> (A B C)
\end{lstlisting}


\small\color{black!55}\hyperref[sol:46]{$\rightarrow$ Soluzione}\normalsize\color{black}
\end{colbox}

\label{ex:47}
\begin{colbox}{47}{Prefisso}{2}
\exlabel{Implementare} \fn{(prefisso-p sub lst)}.
\nopagebreak
\begin{lstlisting}
(prefisso-p '(a b) '(a b c d))   ; -> T
(prefisso-p '(a c) '(a b c d))   ; -> NIL
(prefisso-p '() '(a b))          ; -> T
\end{lstlisting}


\small\color{black!55}\hyperref[sol:47]{$\rightarrow$ Soluzione}\normalsize\color{black}
\end{colbox}

\label{ex:48}
\begin{colbox}{48}{Intercala}{2}
\exlabel{Implementare} \fn{(intercala lst1 lst2)} che alterna gli elementi.
\nopagebreak
\begin{lstlisting}
(intercala '(1 2 3) '(a b c))   ; -> (1 A 2 B 3 C)
\end{lstlisting}
Assumere lunghezze uguali.


\small\color{black!55}\hyperref[sol:48]{$\rightarrow$ Soluzione}\normalsize\color{black}
\end{colbox}

\label{ex:49}
\begin{stdbox}{49}{Run-length encoding}{3}
\textbf{Parte A.} Implementare \fn{(rle-encode lst)}.
\begin{lstlisting}
(rle-encode '(a a a b b c))   ; -> ((3 A) (2 B) (1 C))
\end{lstlisting}
\textbf{Parte B.} Implementare \fn{(rle-decode lst)}: l'inverso.
\begin{lstlisting}
(rle-decode '((3 A) (2 B) (1 C)))   ; -> (A A A B B C)
\end{lstlisting}


\small\color{black!55}\hyperref[sol:49]{$\rightarrow$ Soluzione}\enspace\textcolor{black!30}{|}\enspace\hyperref[hint:49]{$\rightarrow$ Suggerimento}\normalsize\color{black}
\end{stdbox}

\label{ex:50}
\begin{colbox}{50}{Rimozione duplicati}{3}
\exlabel{Implementare} \fn{(remove-duplicates* lst)}: mantiene la prima
occorrenza, senza usare \fn{remove-duplicates}.
\begin{lstlisting}
(remove-duplicates* '(a b a c b a))   ; -> (A B C)
\end{lstlisting}


\small\color{black!55}\hyperref[sol:50]{$\rightarrow$ Soluzione}\enspace\textcolor{black!30}{|}\enspace\hyperref[hint:50]{$\rightarrow$ Suggerimento}\normalsize\color{black}
\end{colbox}

% =============================================================
\setcounter{section}{2}
\section{Ricorsione avanzata}
% =============================================================

\noindent
Tail recursion con accumulatori espliciti, ricorsione su strutture
annidate (alberi impliciti), esponenziazione per quadratura
e primo uso degli array con il crivello di Eratostene.

\bigskip

\label{ex:51}
\begin{essbox}{51}{Fattoriale tail-recursive}{2}
\exlabel{Implementare} \fn{(fattoriale n)} in forma tail-recursive
con accumulatore, usando \fn{labels}.
\begin{lstlisting}
(fattoriale 0)   ; -> 1
(fattoriale 5)   ; -> 120
\end{lstlisting}
\exnota{Obiettivo: capire il pattern accumulatore.}


\small\color{black!55}\hyperref[sol:51]{$\rightarrow$ Soluzione}\normalsize\color{black}
\end{essbox}

\label{ex:52}
\begin{stdbox}{52}{Reverse tail-recursive}{2}
Riscrivere \fn{reverse*} (es.~39) in forma tail-recursive.


\small\color{black!55}\hyperref[sol:52]{$\rightarrow$ Soluzione}\normalsize\color{black}
\end{stdbox}

\label{ex:53}
\begin{colbox}{53}{Somma di un intervallo}{2}
\exlabel{Implementare} \fn{(somma-range a b)} in forma tail-recursive.
\nopagebreak
\begin{lstlisting}
(somma-range 1 100)   ; -> 5050
\end{lstlisting}


\small\color{black!55}\hyperref[sol:53]{$\rightarrow$ Soluzione}\normalsize\color{black}
\end{colbox}

\label{ex:54}
\begin{stdbox}{54}{Fibonacci O(n)}{3}
\exlabel{Implementare} \fn{(fibonacci-fast n)} con due accumulatori,
complessita' $O(n)$.
\begin{lstlisting}
(fibonacci-fast 10)   ; -> 55
(fibonacci-fast 50)   ; -> 12586269025
\end{lstlisting}


\small\color{black!55}\hyperref[sol:54]{$\rightarrow$ Soluzione}\enspace\textcolor{black!30}{|}\enspace\hyperref[hint:54]{$\rightarrow$ Suggerimento}\normalsize\color{black}
\end{stdbox}
\label{ex:55}
\begin{essbox}{55}{Flatten}{3}
\exlabel{Implementare} \fn{(flatten* tree)} che appiattisce strutture annidate.
\nopagebreak
\begin{lstlisting}
(flatten* '(a (b (c d) e) f))   ; -> (A B C D E F)
(flatten* '((1 2) (3 (4 5))))   ; -> (1 2 3 4 5)
\end{lstlisting}


\small\color{black!55}\hyperref[sol:55]{$\rightarrow$ Soluzione}\enspace\textcolor{black!30}{|}\enspace\hyperref[hint:55]{$\rightarrow$ Suggerimento}\normalsize\color{black}
\end{essbox}

\label{ex:56}
\begin{stdbox}{56}{Profondita' di una struttura}{3}
\exlabel{Implementare} \fn{(tree-depth tree)}: profondita' massima di nesting.
\nopagebreak
\begin{lstlisting}
(tree-depth '())             ; -> 0
(tree-depth '(a b c))        ; -> 1
(tree-depth '(a (b c)))      ; -> 2
(tree-depth '(a (b (c d))))  ; -> 3
\end{lstlisting}


\small\color{black!55}\hyperref[sol:56]{$\rightarrow$ Soluzione}\enspace\textcolor{black!30}{|}\enspace\hyperref[hint:56]{$\rightarrow$ Suggerimento}\normalsize\color{black}
\end{stdbox}
\label{ex:57}
\begin{colbox}{57}{Conta atomi}{2}
\exlabel{Implementare} \fn{(count-atoms tree)}.
\nopagebreak
\begin{lstlisting}
(count-atoms '(a (b c) ((d))))   ; -> 4
\end{lstlisting}


\small\color{black!55}\hyperref[sol:57]{$\rightarrow$ Soluzione}\normalsize\color{black}
\end{colbox}

\label{ex:58}
\begin{colbox}{58}{Somma profonda}{2}
\exlabel{Implementare} \fn{(sum-tree tree)}.
\nopagebreak
\begin{lstlisting}
(sum-tree '(1 (2 3) ((4 5))))   ; -> 15
\end{lstlisting}


\small\color{black!55}\hyperref[sol:58]{$\rightarrow$ Soluzione}\normalsize\color{black}
\end{colbox}

\label{ex:59}
\begin{colbox}{59}{Ricerca in struttura annidata}{3}
\exlabel{Implementare} \fn{(tree-member-p item tree)}.
\nopagebreak
\begin{lstlisting}
(tree-member-p 'c '(a (b c) d))   ; -> T
(tree-member-p 'x '(a (b c) d))   ; -> NIL
\end{lstlisting}


\small\color{black!55}\hyperref[sol:59]{$\rightarrow$ Soluzione}\normalsize\color{black}
\end{colbox}

\label{ex:60}
\begin{stdbox}{60}{Sostituzione profonda}{3}
\exlabel{Implementare} \fn{(sostituisci-profondo old new tree)}.
\nopagebreak
\begin{lstlisting}
(sostituisci-profondo 'b 'x '(a (b (b c)) b))
; -> (A (X (X C)) X)
\end{lstlisting}


\small\color{black!55}\hyperref[sol:60]{$\rightarrow$ Soluzione}\enspace\textcolor{black!30}{|}\enspace\hyperref[hint:60]{$\rightarrow$ Suggerimento}\normalsize\color{black}
\end{stdbox}
\label{ex:61}
\begin{stdbox}{61}{Potenza efficiente}{3}
\exlabel{Implementare} \fn{(power-fast base exp)} con complessita' $O(\log n)$:
\[
x^n = \begin{cases}
  (x^{n/2})^2 & \text{se } n \text{ pari} \\
  x \cdot x^{n-1} & \text{se } n \text{ dispari}
\end{cases}
\]
\begin{lstlisting}
(power-fast 2 10)   ; -> 1024
(power-fast 3  0)   ; -> 1
\end{lstlisting}


\small\color{black!55}\hyperref[sol:61]{$\rightarrow$ Soluzione}\enspace\textcolor{black!30}{|}\enspace\hyperref[hint:61]{$\rightarrow$ Suggerimento}\normalsize\color{black}
\end{stdbox}
\label{ex:62}
\begin{stdbox}{62}{MCD e MCM}{2}
Implementare \fn{(mcd* a b)} e \fn{(mcm* a b)} senza usare
\fn{gcd}/\fn{lcm}.
\begin{lstlisting}
(mcd* 48 18)   ; -> 6
(mcm* 4  6)    ; -> 12
\end{lstlisting}


\small\color{black!55}\hyperref[sol:62]{$\rightarrow$ Soluzione}\normalsize\color{black}
\end{stdbox}
\label{ex:63}
\begin{colbox}{63}{Numero primo}{2}
\exlabel{Implementare} \fn{(primo-p n)}.
\nopagebreak
\begin{lstlisting}
(primo-p 17)   ; -> T
(primo-p 18)   ; -> NIL
\end{lstlisting}
Gestire correttamente 0 e 1.


\small\color{black!55}\hyperref[sol:63]{$\rightarrow$ Soluzione}\normalsize\color{black}
\end{colbox}

\label{ex:64}
\begin{essbox}{64}{Crivello di Eratostene}{3}
\exlabel{Implementare} \fn{(crivello n)}: lista di tutti i primi $\leq n$.
\nopagebreak
\begin{lstlisting}
(crivello 20)
; -> (2 3 5 7 11 13 17 19)
\end{lstlisting}
Vincoli: usare array booleano di lunghezza $n+1$; marcare i
multipli con \fn{setf}/\fn{aref}.



\small\color{black!55}\hyperref[sol:64]{$\rightarrow$ Soluzione}\enspace\textcolor{black!30}{|}\enspace\hyperref[hint:64]{$\rightarrow$ Suggerimento}\normalsize\color{black}
\end{essbox}

\label{ex:65}
\begin{stdbox}{65}{Fattori primi}{3}
\exlabel{Implementare} \fn{(prime-factors n)}.
\nopagebreak
\begin{lstlisting}
(prime-factors 60)   ; -> (2 2 3 5)
\end{lstlisting}
Risultato ordinato; $n > 1$.



\small\color{black!55}\hyperref[sol:65]{$\rightarrow$ Soluzione}\enspace\textcolor{black!30}{|}\enspace\hyperref[hint:65]{$\rightarrow$ Suggerimento}\normalsize\color{black}
\end{stdbox}

\label{ex:66}
\begin{colbox}{66}{Sottolista per posizione}{2}
\exlabel{Implementare} \fn{(sublist lst inizio fine)} (indici 0-based,
\fn{fine} escluso).
\begin{lstlisting}
(sublist '(a b c d e) 1 4)   ; -> (B C D)
\end{lstlisting}


\small\color{black!55}\hyperref[sol:66]{$\rightarrow$ Soluzione}\normalsize\color{black}
\end{colbox}

\label{ex:67}
\begin{colbox}{67}{Raggruppa in blocchi}{3}
\exlabel{Implementare} \fn{(raggruppa lst n)}: sottoliste di lunghezza \fn{n}.
\nopagebreak
\begin{lstlisting}
(raggruppa '(a b c d e f g) 3)
; -> ((A B C) (D E F) (G))
\end{lstlisting}


\small\color{black!55}\hyperref[sol:67]{$\rightarrow$ Soluzione}\enspace\textcolor{black!30}{|}\enspace\hyperref[hint:67]{$\rightarrow$ Suggerimento}\normalsize\color{black}
\end{colbox}

\label{ex:68}
\begin{colbox}{68}{Intersezione di insiemi}{2}
\exlabel{Implementare} \fn{(interseca lst1 lst2)} senza usare \fn{intersection}.
\nopagebreak
\begin{lstlisting}
(interseca '(a b c d) '(b d e))   ; -> (B D)
\end{lstlisting}


\small\color{black!55}\hyperref[sol:68]{$\rightarrow$ Soluzione}\normalsize\color{black}
\end{colbox}

% =============================================================
\setcounter{section}{3}
\section{Higher-order e closure}
% =============================================================

\noindent
Funzioni di ordine superiore (\fn{mapcar}, \fn{remove-if},
\fn{reduce}, \fn{apply}, \fn{mapcan}) e closure: funzioni che
catturano variabili dall'ambiente lessicale.

\bigskip

\label{ex:69}
\begin{essbox}{69}{Quadrato degli elementi con mapcar}{1}
\exlabel{Implementare} \fn{(square-list lst)} usando \fn{mapcar} e \fn{lambda}.
\nopagebreak
\begin{lstlisting}
(square-list '(1 2 3 4))   ; -> (1 4 9 16)
\end{lstlisting}


\small\color{black!55}\hyperref[sol:69]{$\rightarrow$ Soluzione}\normalsize\color{black}
\end{essbox}

\label{ex:70}
\begin{stdbox}{70}{Filtro con remove-if-not}{1}
\exlabel{Implementare} \fn{(even-list lst)} usando \fn{remove-if-not}.
\nopagebreak
\begin{lstlisting}
(even-list '(1 2 3 4 5 6))   ; -> (2 4 6)
\end{lstlisting}


\small\color{black!55}\hyperref[sol:70]{$\rightarrow$ Soluzione}\normalsize\color{black}
\end{stdbox}

\label{ex:71}
\begin{colbox}{71}{Stringhe che iniziano con vocale}{2}
\exlabel{Implementare} \fn{(vocali-iniziali lst)} usando \fn{remove-if-not}.
\nopagebreak
\begin{lstlisting}
(vocali-iniziali '("anna" "bob" "elena" "carlo"))
; -> ("anna" "elena")
\end{lstlisting}


\small\color{black!55}\hyperref[sol:71]{$\rightarrow$ Soluzione}\normalsize\color{black}
\end{colbox}

\label{ex:72}
\begin{essbox}{72}{Somma con reduce}{2}
\exlabel{Implementare} \fn{(sum-reduce lst)} con \fn{reduce}.
\nopagebreak
\begin{lstlisting}
(sum-reduce '(1 2 3 4))   ; -> 10
(sum-reduce '())          ; -> 0
\end{lstlisting}


\small\color{black!55}\hyperref[sol:72]{$\rightarrow$ Soluzione}\normalsize\color{black}
\end{essbox}

\label{ex:73}
\begin{colbox}{73}{Prodotto con reduce}{2}
\exlabel{Implementare} \fn{(product-reduce lst)} con \fn{reduce}.
\nopagebreak
\begin{lstlisting}
(product-reduce '(2 3 4))   ; -> 24
\end{lstlisting}


\small\color{black!55}\hyperref[sol:73]{$\rightarrow$ Soluzione}\normalsize\color{black}
\end{colbox}

\label{ex:74}
\begin{colbox}{74}{Min e max con reduce}{2}
\exlabel{Implementare} \fn{(reduce-min lst)} e \fn{(reduce-max lst)}.
\nopagebreak
\begin{lstlisting}
(reduce-min '(8 3 5 1 9))   ; -> 1
(reduce-max '(8 3 5 1 9))   ; -> 9
\end{lstlisting}
Gestire la lista vuota.


\small\color{black!55}\hyperref[sol:74]{$\rightarrow$ Soluzione}\normalsize\color{black}
\end{colbox}

\label{ex:75}
\begin{stdbox}{75}{Lista di cifre in numero intero}{3}
\exlabel{Implementare} \fn{(cifre->intero lst)} con \fn{reduce}.
\nopagebreak
\begin{lstlisting}
(cifre->intero '(1 5 3 4))   ; -> 1534
\end{lstlisting}
Suggerimento: a ogni passo moltiplicare l'accumulatore per 10
e sommare la cifra corrente.



\small\color{black!55}\hyperref[sol:75]{$\rightarrow$ Soluzione}\enspace\textcolor{black!30}{|}\enspace\hyperref[hint:75]{$\rightarrow$ Suggerimento}\normalsize\color{black}
\end{stdbox}

\label{ex:76}
\begin{colbox}{76}{Appiattisce un livello}{3}
\exlabel{Implementare} \fn{(appiattisci-un-livello lst)} con \fn{reduce}
e \fn{append}.
\begin{lstlisting}
(appiattisci-un-livello '((a b) (c d) (e)))
; -> (A B C D E)
\end{lstlisting}


\small\color{black!55}\hyperref[sol:76]{$\rightarrow$ Soluzione}\normalsize\color{black}
\end{colbox}

\label{ex:77}
\begin{colbox}{77}{Regola di Horner}{3}
\exlabel{Implementare} \fn{(horner coefficienti x)} con \fn{reduce}.
\[
p(x) = (\ldots((c_n \cdot x + c_{n-1}) \cdot x + c_{n-2})\ldots + c_0)
\]
\begin{lstlisting}
; 2x^2 + 3x + 1  con x=2  -> 15
(horner '(2 3 1) 2)   ; -> 15
\end{lstlisting}


\small\color{black!55}\hyperref[sol:77]{$\rightarrow$ Soluzione}\enspace\textcolor{black!30}{|}\enspace\hyperref[hint:77]{$\rightarrow$ Suggerimento}\normalsize\color{black}
\end{colbox}

\label{ex:78}
\begin{stdbox}{78}{Prodotto scalare con apply}{2}
\exlabel{Implementare} \fn{(prodotto-scalare v1 v2)}.
\nopagebreak
\begin{lstlisting}
(prodotto-scalare '(1 2 3) '(4 5 6))   ; -> 32
\end{lstlisting}
Vincolo: usare \fn{mapcar} per i prodotti componente per componente,
poi \fn{apply} con \fn{+} per la somma.


\small\color{black!55}\hyperref[sol:78]{$\rightarrow$ Soluzione}\normalsize\color{black}
\end{stdbox}

\label{ex:79}
\begin{stdbox}{79}{mapcan}{2}
\exlabel{Implementare} \fn{(raddoppia lst)} con \fn{mapcan}.
\nopagebreak
\begin{lstlisting}
(raddoppia '(1 2 3))   ; -> (1 1 2 2 3 3)
\end{lstlisting}
Spiegare la differenza tra \fn{mapcan} e \fn{mapcar} in un commento.


\small\color{black!55}\hyperref[sol:79]{$\rightarrow$ Soluzione}\normalsize\color{black}
\end{stdbox}

\label{ex:80}
\begin{colbox}{80}{Ordina per chiave}{2}
\exlabel{Implementare} \fn{(ordina-per-punteggio lst)}: ordina coppie
\fn{(nome . punteggio)} in ordine decrescente.
\begin{lstlisting}
(ordina-per-punteggio
  '((mario . 85) (anna . 92) (luca . 78)))
; -> ((ANNA . 92) (MARIO . 85) (LUCA . 78))
\end{lstlisting}
Vincolo: usare \fn{sort} con comparatore custom.


\small\color{black!55}\hyperref[sol:80]{$\rightarrow$ Soluzione}\normalsize\color{black}
\end{colbox}

\label{ex:81}
\begin{colbox}{81}{Tutti e qualcuno}{2}
\exlabel{Implementare} \fn{(tutti-p pred lst)} e \fn{(qualcuno-p pred lst)}
replicando \fn{every}/\fn{some} senza usarli.
\begin{lstlisting}
(tutti-p #'evenp '(2 4 6))    ; -> T
(qualcuno-p #'oddp '(2 3 4))  ; -> T
\end{lstlisting}


\small\color{black!55}\hyperref[sol:81]{$\rightarrow$ Soluzione}\normalsize\color{black}
\end{colbox}

\label{ex:82}
\begin{colbox}{82}{Conta con predicato}{2}
\exlabel{Implementare} \fn{(count-matching pred lst)} senza usare \fn{count-if}.
\nopagebreak
\begin{lstlisting}
(count-matching #'evenp '(1 2 3 4 5))   ; -> 2
\end{lstlisting}


\small\color{black!55}\hyperref[sol:82]{$\rightarrow$ Soluzione}\normalsize\color{black}
\end{colbox}

\label{ex:83}
\begin{colbox}{83}{Raggruppa per chiave}{3}
\exlabel{Implementare} \fn{(raggruppa-per funzione lst)}: restituisce
una alist con chiave = risultato della funzione.
\begin{lstlisting}
(raggruppa-per #'evenp '(1 2 3 4 5))
; -> ((NIL 1 3 5) (T 2 4))
\end{lstlisting}


\small\color{black!55}\hyperref[sol:83]{$\rightarrow$ Soluzione}\enspace\textcolor{black!30}{|}\enspace\hyperref[hint:83]{$\rightarrow$ Suggerimento}\normalsize\color{black}
\end{colbox}

\label{ex:84}
\begin{colbox}{84}{Applica a tutti}{2}
\exlabel{Implementare} \fn{(applica-a-tutti funzioni x)}.
\nopagebreak
\begin{lstlisting}
(applica-a-tutti (list #'1+ #'1- #'zerop) 0)
; -> (1 -1 T)
\end{lstlisting}


\small\color{black!55}\hyperref[sol:84]{$\rightarrow$ Soluzione}\normalsize\color{black}
\end{colbox}

\label{ex:85}
\begin{essbox}{85}{Composizione di funzioni}{3}
\exlabel{Implementare} \fn{(compose f g)}: restituisce $f(g(x))$ come closure.
\nopagebreak
\begin{lstlisting}
(defparameter *f* (compose #'1+ #'1+))
(funcall *f* 10)   ; -> 12
\end{lstlisting}


\small\color{black!55}\hyperref[sol:85]{$\rightarrow$ Soluzione}\enspace\textcolor{black!30}{|}\enspace\hyperref[hint:85]{$\rightarrow$ Suggerimento}\normalsize\color{black}
\end{essbox}

\label{ex:86}
\begin{stdbox}{86}{Applicazione parziale}{3}
\exlabel{Implementare} \fn{(partial funzione \&rest argomenti)}.
\nopagebreak
\begin{lstlisting}
(defparameter *add10* (partial #'+ 10))
(funcall *add10* 5)    ; -> 15
(funcall *add10* 20)   ; -> 30
\end{lstlisting}


\small\color{black!55}\hyperref[sol:86]{$\rightarrow$ Soluzione}\enspace\textcolor{black!30}{|}\enspace\hyperref[hint:86]{$\rightarrow$ Suggerimento}\normalsize\color{black}
\end{stdbox}

\label{ex:87}
\begin{essbox}{87}{Closure con stato: contatore}{3}
\exlabel{Implementare} \fn{(make-counter \&optional (start 0) (step 1))}.
\nopagebreak
\begin{lstlisting}
(defparameter *c* (make-counter))
(funcall *c*)   ; -> 0
(funcall *c*)   ; -> 1

(defparameter *c2* (make-counter 10 5))
(funcall *c2*)  ; -> 10
(funcall *c2*)  ; -> 15
\end{lstlisting}


\small\color{black!55}\hyperref[sol:87]{$\rightarrow$ Soluzione}\enspace\textcolor{black!30}{|}\enspace\hyperref[hint:87]{$\rightarrow$ Suggerimento}\normalsize\color{black}
\end{essbox}

\label{ex:88}
\begin{colbox}{88}{Accumulatore}{3}
\exlabel{Implementare} \fn{(make-accumulator iniziale)}: ogni chiamata
somma un valore al totale.
\begin{lstlisting}
(defparameter *acc* (make-accumulator 5))
(funcall *acc* 10)   ; -> 15
(funcall *acc* 3)    ; -> 18
\end{lstlisting}


\small\color{black!55}\hyperref[sol:88]{$\rightarrow$ Soluzione}\normalsize\color{black}
\end{colbox}

\label{ex:89}
\begin{colbox}{89}{Generatore Fibonacci con closure}{4}
\exlabel{Implementare} \fn{(make-fib-generator)}: restituisce una funzione
che ad ogni chiamata produce il successivo numero di Fibonacci.
\begin{lstlisting}
(defparameter *fib* (make-fib-generator))
(funcall *fib*)   ; -> 0
(funcall *fib*)   ; -> 1
(funcall *fib*)   ; -> 1
(funcall *fib*)   ; -> 2
\end{lstlisting}


\small\color{black!55}\hyperref[sol:89]{$\rightarrow$ Soluzione}\enspace\textcolor{black!30}{|}\enspace\hyperref[hint:89]{$\rightarrow$ Suggerimento}\normalsize\color{black}
\end{colbox}

\label{ex:90}
\begin{colbox}{90}{Pipeline}{3}
\exlabel{Implementare} \fn{(pipeline valore funzioni)}.
\nopagebreak
\begin{lstlisting}
(pipeline 5 (list #'1+ #'1+ (lambda (x) (* x 2))))
; -> 14
\end{lstlisting}


\small\color{black!55}\hyperref[sol:90]{$\rightarrow$ Soluzione}\normalsize\color{black}
\end{colbox}

\label{ex:91}
\begin{stdbox}{91}{Memoization}{4}
\exlabel{Implementare} \fn{(memoize funzione)} che restituisce una versione
con cache.

Vincoli: hash table nella closure; funzione originale invariata.
\begin{lstlisting}
(defparameter *fib-m*
  (memoize #'fibonacci))
(funcall *fib-m* 35)   ; veloce
\end{lstlisting}


\small\color{black!55}\hyperref[sol:91]{$\rightarrow$ Soluzione}\enspace\textcolor{black!30}{|}\enspace\hyperref[hint:91]{$\rightarrow$ Suggerimento}\normalsize\color{black}
\end{stdbox}

\label{ex:92}
\begin{colbox}{92}{Mini libreria funzionale}{4}
Raccogliere in un unico file, con docstring:
\fn{map-list}, \fn{filter-list}, \fn{count-matching},
\fn{find-matching}, \fn{compose}, \fn{partial}.

Aggiungere almeno 3 funzioni a scelta (es.\ \fn{take}, \fn{drop}, \fn{zip}).

Requisiti: nessuno stato globale; almeno un test per funzione.


\small\color{black!55}\hyperref[sol:92]{$\rightarrow$ Soluzione}\normalsize\color{black}
\end{colbox}

% =============================================================
\setcounter{section}{4}
\section{Strutture dati}
% =============================================================

\noindent
Association list, hash table, array 2D, \fn{defstruct},
alberi binari e BST, pila e coda come ADT, grafi.

\bigskip

\label{ex:93}
\begin{essbox}{93}{Association list}{2}
\exlabel{Implementare} \fn{(lookup key alist)} e \fn{(set-key key value alist)}.
\nopagebreak
\begin{lstlisting}
(lookup 'age '((name . "Mario") (age . 23)))
; -> 23
(set-key 'age 24 '((name . "Mario") (age . 23)))
; -> ((NAME . "Mario") (AGE . 24))
\end{lstlisting}
Usare \fn{assoc} in \fn{lookup}; non modificare la lista originale.


\small\color{black!55}\hyperref[sol:93]{$\rightarrow$ Soluzione}\normalsize\color{black}
\end{essbox}

\label{ex:94}
\begin{stdbox}{94}{Frequenza elementi}{2}
\exlabel{Implementare} \fn{(frequencies lst)}.
\nopagebreak
\begin{lstlisting}
(frequencies '(a b a c b a))
; -> ((A . 3) (B . 2) (C . 1))
\end{lstlisting}


\small\color{black!55}\hyperref[sol:94]{$\rightarrow$ Soluzione}\enspace\textcolor{black!30}{|}\enspace\hyperref[hint:94]{$\rightarrow$ Suggerimento}\normalsize\color{black}
\end{stdbox}

\label{ex:95}
\begin{essbox}{95}{Hash table: operazioni base}{2}
\exlabel{Implementare} \fn{(ht-get key ht)}, \fn{(ht-set key value ht)},
\fn{(ht-delete key ht)}, \fn{(ht-keys ht)}.
\begin{lstlisting}
(defparameter *ht* (make-hash-table))
(ht-set 'foo 42 *ht*)
(ht-get 'foo *ht*)     ; -> 42
(ht-delete 'foo *ht*)
(ht-get 'foo *ht*)     ; -> NIL
\end{lstlisting}


\small\color{black!55}\hyperref[sol:95]{$\rightarrow$ Soluzione}\normalsize\color{black}
\end{essbox}

\label{ex:96}
\begin{colbox}{96}{Frequenza con hash table}{2}
Riscrivere \fn{(frequencies-hash lst)} usando una hash table.

\exnota{Obiettivo: confrontare alist vs hash table.}


\small\color{black!55}\hyperref[sol:96]{$\rightarrow$ Soluzione}\normalsize\color{black}
\end{colbox}

\label{ex:97}
\begin{stdbox}{97}{Frequenza dei caratteri}{3}
\exlabel{Implementare} \fn{(frequenza-caratteri stringa)}.
\nopagebreak
\begin{lstlisting}
(frequenza-caratteri "banana")
; hash table: #\b->1, #\a->3, #\n->2
\end{lstlisting}


\small\color{black!55}\hyperref[sol:97]{$\rightarrow$ Soluzione}\enspace\textcolor{black!30}{|}\enspace\hyperref[hint:97]{$\rightarrow$ Suggerimento}\normalsize\color{black}
\end{stdbox}

\label{ex:98}
\begin{colbox}{98}{Unione di hash table}{3}
\exlabel{Implementare} \fn{(unisci-hash ht1 ht2)}: chiavi duplicate
$\to$ somma dei valori.



\small\color{black!55}\hyperref[sol:98]{$\rightarrow$ Soluzione}\enspace\textcolor{black!30}{|}\enspace\hyperref[hint:98]{$\rightarrow$ Suggerimento}\normalsize\color{black}
\end{colbox}

\label{ex:99}
\begin{colbox}{99}{Inverti hash table}{3}
\exlabel{Implementare} \fn{(inverti-hash ht)}: scambia chiavi e valori.
Assumere valori unici.



\small\color{black!55}\hyperref[sol:99]{$\rightarrow$ Soluzione}\enspace\textcolor{black!30}{|}\enspace\hyperref[hint:99]{$\rightarrow$ Suggerimento}\normalsize\color{black}
\end{colbox}

\label{ex:100}
\begin{stdbox}{100}{Conversione alist/hash}{2}
\exlabel{Implementare} \fn{(alist->hash alist)} e \fn{(hash->alist ht)}.
\nopagebreak
\begin{lstlisting}
(alist->hash '((a . 1) (b . 2)))
; -> hash table a->1, b->2
\end{lstlisting}


\small\color{black!55}\hyperref[sol:100]{$\rightarrow$ Soluzione}\normalsize\color{black}
\end{stdbox}

\label{ex:101}
\begin{essbox}{101}{Matrice con array 2D}{2}
Creare una matrice $3 \times 3$ con \fn{make-array} e implementare:
\begin{lstlisting}
(mat-get m r c)
(mat-set m r c val)
(stampa-matrice m)
\end{lstlisting}
\begin{lstlisting}
(defparameter *m* (make-array '(3 3) :initial-element 0))
(mat-set *m* 1 1 5)
(mat-get *m* 1 1)   ; -> 5
\end{lstlisting}


\small\color{black!55}\hyperref[sol:101]{$\rightarrow$ Soluzione}\normalsize\color{black}
\end{essbox}

\label{ex:102}
\begin{colbox}{102}{Trasposta di matrice}{3}
\exlabel{Implementare} \fn{(trasposta m n)}: trasposta di una matrice $n \times n$.



\small\color{black!55}\hyperref[sol:102]{$\rightarrow$ Soluzione}\enspace\textcolor{black!30}{|}\enspace\hyperref[hint:102]{$\rightarrow$ Suggerimento}\normalsize\color{black}
\end{colbox}

\label{ex:103}
\begin{stdbox}{103}{Inversione array in loco}{2}
\exlabel{Implementare} \fn{(inverti-array! v)} che inverte il vettore
modificandolo con \fn{setf}/\fn{aref}, senza liste ausiliarie.
\begin{lstlisting}
(defparameter *v*
  (make-array 5 :initial-contents '(1 2 3 4 5)))
(inverti-array! *v*)
; *v* ora e': #(5 4 3 2 1)
\end{lstlisting}


\small\color{black!55}\hyperref[sol:103]{$\rightarrow$ Soluzione}\normalsize\color{black}
\end{stdbox}

\label{ex:104}
\begin{colbox}{104}{Moltiplicazione di matrici (opzionale)}{4}
\exlabel{Implementare} \fn{(moltiplica-matrici a b n)}: prodotto di
matrici quadrate $n \times n$.



\small\color{black!55}\hyperref[sol:104]{$\rightarrow$ Soluzione}\enspace\textcolor{black!30}{|}\enspace\hyperref[hint:104]{$\rightarrow$ Suggerimento}\normalsize\color{black}
\end{colbox}

\label{ex:105}
\begin{stdbox}{105}{Struttura con defstruct}{2}
Definire:
\begin{lstlisting}
(defstruct persona
  nome
  (eta 18)
  citta)
\end{lstlisting}
\exlabel{Implementare} \fn{(crea-persona nome eta citta)} che valida l'eta'
(positiva) e segnala errore altrimenti.
\begin{lstlisting}
(crea-persona "Mario" 23 "Roma")
; -> #S(PERSONA :NOME "Mario" :ETA 23 :CITTA "Roma")
\end{lstlisting}


\small\color{black!55}\hyperref[sol:105]{$\rightarrow$ Soluzione}\normalsize\color{black}
\end{stdbox}

\label{ex:106}
\begin{colbox}{106}{Struttura nodo-albero}{2}
Definire:
\begin{lstlisting}
(defstruct nodo
  valore
  (sinistra nil)
  (destra   nil))
\end{lstlisting}
Costruire manualmente l'albero:
\begin{lstlisting}
;       5
;      / \
;     3   8
;    / \
;   1   4
\end{lstlisting}


\small\color{black!55}\hyperref[sol:106]{$\rightarrow$ Soluzione}\normalsize\color{black}
\end{colbox}

\label{ex:107}
\begin{colbox}{107}{Albero con liste}{2}
Rappresentare lo stesso albero con la notazione
\fn{(valore sinistra destra)} e implementare:
\begin{lstlisting}
(make-leaf v)
(make-node v left right)
\end{lstlisting}


\small\color{black!55}\hyperref[sol:107]{$\rightarrow$ Soluzione}\normalsize\color{black}
\end{colbox}

\label{ex:108}
\begin{stdbox}{108}{BST: inserimento e ricerca}{4}
Implementare un albero binario di ricerca:
\begin{lstlisting}
(bst-insert value tree)
(bst-find value tree)
\end{lstlisting}
\begin{lstlisting}
(bst-find 4 (bst-insert 4 (bst-insert 5 nil)))
; -> T
\end{lstlisting}


\small\color{black!55}\hyperref[sol:108]{$\rightarrow$ Soluzione}\enspace\textcolor{black!30}{|}\enspace\hyperref[hint:108]{$\rightarrow$ Suggerimento}\normalsize\color{black}
\end{stdbox}

\label{ex:109}
\begin{stdbox}{109}{Traversal dell'albero}{3}
\exlabel{Implementare} \fn{(preorder tree)}, \fn{(inorder tree)},
\fn{(postorder tree)}.

Verificare che \fn{inorder} su un BST produca sequenza ordinata.
\begin{lstlisting}
(inorder albero)   ; -> (1 3 4 5 8)
\end{lstlisting}


\small\color{black!55}\hyperref[sol:109]{$\rightarrow$ Soluzione}\enspace\textcolor{black!30}{|}\enspace\hyperref[hint:109]{$\rightarrow$ Suggerimento}\normalsize\color{black}
\end{stdbox}

\label{ex:110}
\begin{colbox}{110}{Altezza e conteggio foglie}{3}
\exlabel{Implementare} \fn{(tree-height tree)} e \fn{(count-leaves tree)}.
\nopagebreak
\begin{lstlisting}
(tree-height albero)    ; -> 3
(count-leaves albero)   ; -> 3
\end{lstlisting}


\small\color{black!55}\hyperref[sol:110]{$\rightarrow$ Soluzione}\normalsize\color{black}
\end{colbox}

\label{ex:111}
\begin{colbox}{111}{Specchia l'albero}{3}
\exlabel{Implementare} \fn{(specchia tree)}: scambia ricorsivamente
sottoalberi sinistro e destro.


\small\color{black!55}\hyperref[sol:111]{$\rightarrow$ Soluzione}\normalsize\color{black}
\end{colbox}

\label{ex:112}
\begin{colbox}{112}{Albero bilanciato?}{4}
\exlabel{Implementare} \fn{(bilanciato-p tree)}: le altezze dei sottoalberi
differiscono al piu' di 1 per ogni nodo.



\small\color{black!55}\hyperref[sol:112]{$\rightarrow$ Soluzione}\enspace\textcolor{black!30}{|}\enspace\hyperref[hint:112]{$\rightarrow$ Suggerimento}\normalsize\color{black}
\end{colbox}

\label{ex:113}
\begin{stdbox}{113}{Pila (stack) come ADT}{2}
Implementare con \fn{setf}:
\begin{lstlisting}
(make-stack)
(stack-push item stack)
(stack-pop stack)
(stack-peek stack)
(stack-empty-p stack)
\end{lstlisting}
\begin{lstlisting}
(defparameter *s* (make-stack))
(stack-push 1 *s*)
(stack-push 2 *s*)
(stack-pop *s*)    ; -> 2
(stack-peek *s*)   ; -> 1
\end{lstlisting}


\small\color{black!55}\hyperref[sol:113]{$\rightarrow$ Soluzione}\normalsize\color{black}
\end{stdbox}

\label{ex:114}
\begin{stdbox}{114}{Coda (queue) come ADT}{3}
\exlabel{Implementare:}
\nopagebreak
\begin{lstlisting}
(make-queue)
(enqueue item queue)
(dequeue queue)
(queue-empty-p queue)
\end{lstlisting}


\small\color{black!55}\hyperref[sol:114]{$\rightarrow$ Soluzione}\enspace\textcolor{black!30}{|}\enspace\hyperref[hint:114]{$\rightarrow$ Suggerimento}\normalsize\color{black}
\end{stdbox}

\label{ex:115}
\begin{stdbox}{115}{Valida parentesi}{3}
\exlabel{Implementare} \fn{(valida-parentesi stringa)} per
\fn{()}, \fn{[]}, \fn{\{\}}.
\begin{lstlisting}
(valida-parentesi "([]{})")   ; -> T
(valida-parentesi "([)]")     ; -> NIL
(valida-parentesi "((")       ; -> NIL
\end{lstlisting}
Vincolo: usare una pila.



\small\color{black!55}\hyperref[sol:115]{$\rightarrow$ Soluzione}\enspace\textcolor{black!30}{|}\enspace\hyperref[hint:115]{$\rightarrow$ Suggerimento}\normalsize\color{black}
\end{stdbox}

\label{ex:116}
\begin{colbox}{116}{Rubrica con hash table}{3}
\exlabel{Implementare:}
\nopagebreak
\begin{lstlisting}
(add-contact name phone)
(find-contact name)
(remove-contact name)
(list-contacts)
\end{lstlisting}
Gestire il caso contatto non trovato.


\small\color{black!55}\hyperref[sol:116]{$\rightarrow$ Soluzione}\normalsize\color{black}
\end{colbox}

\label{ex:117}
\begin{stdbox}{117}{Grafo come lista di adiacenza}{3}
\exlabel{Implementare:}
\nopagebreak
\begin{lstlisting}
(neighbors node graph)
(add-edge node1 node2 graph)
\end{lstlisting}
\begin{lstlisting}
(neighbors 'a '((a b c) (b a d) (c a d) (d b c)))
; -> (B C)
\end{lstlisting}


\small\color{black!55}\hyperref[sol:117]{$\rightarrow$ Soluzione}\normalsize\color{black}
\end{stdbox}

\label{ex:118}
\begin{colbox}{118}{DFS su grafo}{4}
\exlabel{Implementare} \fn{(dfs graph start)}: visita in profondita'.
\nopagebreak
\begin{lstlisting}
(dfs '((a b c) (b a d) (c a) (d b)) 'a)
; -> (A B D C)
\end{lstlisting}
Gestire i cicli con un insieme di nodi gia' visitati.



\small\color{black!55}\hyperref[sol:118]{$\rightarrow$ Soluzione}\enspace\textcolor{black!30}{|}\enspace\hyperref[hint:118]{$\rightarrow$ Suggerimento}\normalsize\color{black}
\end{colbox}

\label{ex:119}
\begin{colbox}{119}{BFS su grafo}{4}
\exlabel{Implementare} \fn{(bfs graph start)}: visita in ampiezza.
\nopagebreak
\begin{lstlisting}
(bfs '((a b c) (b a d) (c a) (d b)) 'a)
; -> (A B C D)
\end{lstlisting}
Vincolo: usare una coda.



\small\color{black!55}\hyperref[sol:119]{$\rightarrow$ Soluzione}\enspace\textcolor{black!30}{|}\enspace\hyperref[hint:119]{$\rightarrow$ Suggerimento}\normalsize\color{black}
\end{colbox}

\label{ex:120}
\begin{colbox}{120}{Conta parole}{3}
\exlabel{Implementare} \fn{(conta-parole testo)} e
\fn{(parola-piu-frequente testo)}.
\begin{lstlisting}
(conta-parole "il gatto e' sul tetto e il tetto")
; hash table: "il"->2, ...
\end{lstlisting}


\small\color{black!55}\hyperref[sol:120]{$\rightarrow$ Soluzione}\enspace\textcolor{black!30}{|}\enspace\hyperref[hint:120]{$\rightarrow$ Suggerimento}\normalsize\color{black}
\end{colbox}

% =============================================================
\setcounter{section}{5}
\section{Macro}
% =============================================================

\noindent
Le macro operano a livello sintattico: trasformano codice
prima che venga valutato. Questo capitolo mostra come costruire
nuove strutture di controllo.

\bigskip

\label{ex:121}
\begin{stdbox}{121}{Macro quando}{3}
\exlabel{Implementare:}
\nopagebreak
\begin{lstlisting}
(defmacro quando (test &body corpo))
\end{lstlisting}
Si comporta come \fn{when}: esegue \fn{corpo} solo se \fn{test} e' vero.
\begin{lstlisting}
(quando (> x 0)
  (print "positivo")
  (print "ok"))
\end{lstlisting}
Verificare l'espansione con \fn{macroexpand-1}.

\exnota{Obiettivo: capire la differenza macro/funzione.}



\small\color{black!55}\hyperref[sol:121]{$\rightarrow$ Soluzione}\enspace\textcolor{black!30}{|}\enspace\hyperref[hint:121]{$\rightarrow$ Suggerimento}\normalsize\color{black}
\end{stdbox}

\label{ex:122}
\begin{stdbox}{122}{Macro while}{3}
\exlabel{Implementare:}
\nopagebreak
\begin{lstlisting}
(defmacro mio-while (test &body corpo))
\end{lstlisting}
\begin{lstlisting}
(let ((i 0))
  (mio-while (< i 5)
    (print i)
    (setf i (1+ i))))
; stampa 0 1 2 3 4
\end{lstlisting}


\small\color{black!55}\hyperref[sol:122]{$\rightarrow$ Soluzione}\enspace\textcolor{black!30}{|}\enspace\hyperref[hint:122]{$\rightarrow$ Suggerimento}\normalsize\color{black}
\end{stdbox}

\label{ex:123}
\begin{colbox}{123}{Macro for}{4}
\exlabel{Implementare:}
\nopagebreak
\begin{lstlisting}
(defmacro mio-for ((var da a) &body corpo))
\end{lstlisting}
\begin{lstlisting}
(mio-for (i 1 5)
  (format t "~A~%" i))
; stampa 1 2 3 4 5
\end{lstlisting}
\exnota{Obiettivo: usare \fn{gensym} per evitare cattura della variabile.}



\small\color{black!55}\hyperref[sol:123]{$\rightarrow$ Soluzione}\enspace\textcolor{black!30}{|}\enspace\hyperref[hint:123]{$\rightarrow$ Suggerimento}\normalsize\color{black}
\end{colbox}

% =============================================================
\setcounter{section}{6}
\section{Backtracking e algoritmi}
% =============================================================

\noindent
Algoritmi che esplorano uno spazio di ricerca: generazione di
combinazioni, ordinamento e problemi classici di backtracking.

\bigskip

\label{ex:124}
\begin{stdbox}{124}{Tutti i sottoinsiemi}{3}
\exlabel{Implementare} \fn{(subsets lst)}: genera tutti i $2^n$ sottoinsiemi.
\nopagebreak
\begin{lstlisting}
(subsets '(a b))
; -> (() (A) (B) (A B))
\end{lstlisting}


\small\color{black!55}\hyperref[sol:124]{$\rightarrow$ Soluzione}\enspace\textcolor{black!30}{|}\enspace\hyperref[hint:124]{$\rightarrow$ Suggerimento}\normalsize\color{black}
\end{stdbox}

\label{ex:125}
\begin{stdbox}{125}{Permutazioni}{4}
\exlabel{Implementare} \fn{(permutations lst)}.
\nopagebreak
\begin{lstlisting}
(length (permutations '(a b c)))   ; -> 6
\end{lstlisting}


\small\color{black!55}\hyperref[sol:125]{$\rightarrow$ Soluzione}\enspace\textcolor{black!30}{|}\enspace\hyperref[hint:125]{$\rightarrow$ Suggerimento}\normalsize\color{black}
\end{stdbox}

\label{ex:126}
\begin{stdbox}{126}{Combinazioni}{3}
\exlabel{Implementare} \fn{(combinations lst k)}.
\nopagebreak
\begin{lstlisting}
(combinations '(a b c d) 2)
; -> ((A B) (A C) (A D) (B C) (B D) (C D))
\end{lstlisting}


\small\color{black!55}\hyperref[sol:126]{$\rightarrow$ Soluzione}\enspace\textcolor{black!30}{|}\enspace\hyperref[hint:126]{$\rightarrow$ Suggerimento}\normalsize\color{black}
\end{stdbox}

\label{ex:127}
\begin{stdbox}{127}{Somma a target}{3}
\exlabel{Implementare} \fn{(somma-a-target numeri target)}.
\nopagebreak
\begin{lstlisting}
(somma-a-target '(1 2 3 4 5) 6)
; -> ((1 5) (2 4) (1 2 3))
\end{lstlisting}


\small\color{black!55}\hyperref[sol:127]{$\rightarrow$ Soluzione}\enspace\textcolor{black!30}{|}\enspace\hyperref[hint:127]{$\rightarrow$ Suggerimento}\normalsize\color{black}
\end{stdbox}

\label{ex:128}
\begin{colbox}{128}{Parentesi bilanciate generate}{3}
\exlabel{Implementare} \fn{(genera-parentesi n)}.
\nopagebreak
\begin{lstlisting}
(genera-parentesi 3)
; -> ("((()))" "(()())" "(())()" "()(())" "()()()")
\end{lstlisting}


\small\color{black!55}\hyperref[sol:128]{$\rightarrow$ Soluzione}\enspace\textcolor{black!30}{|}\enspace\hyperref[hint:128]{$\rightarrow$ Suggerimento}\normalsize\color{black}
\end{colbox}

\label{ex:129}
\begin{stdbox}{129}{Merge sort}{4}
\exlabel{Implementare} \fn{(merge-sorted lst1 lst2)} e \fn{(merge-sort* lst)}.
\nopagebreak
\begin{lstlisting}
(merge-sort* '(7 2 9 1 5 3))   ; -> (1 2 3 5 7 9)
\end{lstlisting}
Vincoli: non usare \fn{sort}; complessita' $O(n \log n)$.



\small\color{black!55}\hyperref[sol:129]{$\rightarrow$ Soluzione}\enspace\textcolor{black!30}{|}\enspace\hyperref[hint:129]{$\rightarrow$ Suggerimento}\normalsize\color{black}
\end{stdbox}

\label{ex:130}
\begin{stdbox}{130}{Quicksort}{4}
\exlabel{Implementare} \fn{(quicksort* lst)} senza usare \fn{sort}.
\nopagebreak
\begin{lstlisting}
(quicksort* '(7 2 9 1 5 3))   ; -> (1 2 3 5 7 9)
\end{lstlisting}


\small\color{black!55}\hyperref[sol:130]{$\rightarrow$ Soluzione}\enspace\textcolor{black!30}{|}\enspace\hyperref[hint:130]{$\rightarrow$ Suggerimento}\normalsize\color{black}
\end{stdbox}

\label{ex:131}
\begin{stdbox}{131}{Torri di Hanoi}{3}
\exlabel{Implementare} \fn{(hanoi n da a tramite)} che stampa le mosse.
\nopagebreak
\begin{lstlisting}
(hanoi 3 'A 'C 'B)
; Sposta disco 1: A -> C
; ...
\end{lstlisting}


\small\color{black!55}\hyperref[sol:131]{$\rightarrow$ Soluzione}\enspace\textcolor{black!30}{|}\enspace\hyperref[hint:131]{$\rightarrow$ Suggerimento}\normalsize\color{black}
\end{stdbox}

\label{ex:132}
\begin{colbox}{132}{Problema dello zaino}{4}
\exlabel{Implementare} \fn{(knapsack oggetti capacita)}.
Oggetti: lista di \fn{(peso valore)}.
\begin{lstlisting}
(knapsack '((2 3) (3 4) (4 5) (5 8)) 5)
; -> 7
\end{lstlisting}


\small\color{black!55}\hyperref[sol:132]{$\rightarrow$ Soluzione}\enspace\textcolor{black!30}{|}\enspace\hyperref[hint:132]{$\rightarrow$ Suggerimento}\normalsize\color{black}
\end{colbox}

\label{ex:133}
\begin{colbox}{133}{Labirinto}{4}
Rappresentare un labirinto come matrice 2D (0 libero, 1 muro).
\exlabel{Implementare} \fn{(find-path maze start goal)}.
\nopagebreak
\begin{lstlisting}
(find-path maze '(0 0) '(4 4))
; -> ((0 0) (1 0) ... (4 4))
\end{lstlisting}


\small\color{black!55}\hyperref[sol:133]{$\rightarrow$ Soluzione}\enspace\textcolor{black!30}{|}\enspace\hyperref[hint:133]{$\rightarrow$ Suggerimento}\normalsize\color{black}
\end{colbox}

\label{ex:134}
\begin{stdbox}{134}{N-regine}{5}
\exlabel{Implementare} \fn{(n-queens n)}: tutte le configurazioni valide.
\nopagebreak
\begin{lstlisting}
(length (n-queens 4))   ; -> 2
(length (n-queens 8))   ; -> 92
\end{lstlisting}
Vincoli: backtracking con pruning; rappresentazione compatta.



\small\color{black!55}\hyperref[sol:134]{$\rightarrow$ Soluzione}\enspace\textcolor{black!30}{|}\enspace\hyperref[hint:134]{$\rightarrow$ Suggerimento}\normalsize\color{black}
\end{stdbox}

\label{ex:135}
\begin{colbox}{135}{Sudoku}{5}
\exlabel{Implementare} \fn{(solve-sudoku board)}.
\fn{board}: array $9 \times 9$ con 0 per celle vuote.

Algoritmo: per ogni cella vuota, provare i valori 1--9,
verificare vincoli riga/colonna/blocco, procedere ricorsivamente.



\small\color{black!55}\hyperref[sol:135]{$\rightarrow$ Soluzione}\enspace\textcolor{black!30}{|}\enspace\hyperref[hint:135]{$\rightarrow$ Suggerimento}\normalsize\color{black}
\end{colbox}

\label{ex:136}
\begin{colbox}{136}{Cammino del cavallo}{5}
\exlabel{Implementare} \fn{(knights-tour n)}: cammino del cavallo su
scacchiera $n \times n$ visitando ogni casella esattamente una volta.
\begin{lstlisting}
(knights-tour 5)   ; -> lista di mosse, o NIL
\end{lstlisting}
Suggerimento: euristica di Warnsdorff.



\small\color{black!55}\hyperref[sol:136]{$\rightarrow$ Soluzione}\enspace\textcolor{black!30}{|}\enspace\hyperref[hint:136]{$\rightarrow$ Suggerimento}\normalsize\color{black}
\end{colbox}

% =============================================================
\setcounter{section}{7}
\section{Stringhe e I/O}
% =============================================================

\noindent
Operazioni sulle stringhe e tecniche di I/O.

\bigskip

\label{ex:137}
\begin{stdbox}{137}{Inverti una stringa}{2}
\exlabel{Implementare} \fn{(reverse-string s)} senza applicare \fn{reverse}
direttamente; usare \fn{coerce}.
\begin{lstlisting}
(reverse-string "hello")   ; -> "olleh"
\end{lstlisting}


\small\color{black!55}\hyperref[sol:137]{$\rightarrow$ Soluzione}\normalsize\color{black}
\end{stdbox}

\label{ex:138}
\begin{stdbox}{138}{Stringa palindroma}{2}
\exlabel{Implementare} \fn{(palindrome-p s)}.
\nopagebreak
\begin{lstlisting}
(palindrome-p "anna")    ; -> T
(palindrome-p "casa")    ; -> NIL
\end{lstlisting}


\small\color{black!55}\hyperref[sol:138]{$\rightarrow$ Soluzione}\normalsize\color{black}
\end{stdbox}

\label{ex:139}
\begin{colbox}{139}{Palindromo robusto}{3}
\exlabel{Implementare} \fn{(palindrome-robust-p s)}: ignora
maiuscole/minuscole, spazi, punteggiatura.
\begin{lstlisting}
(palindrome-robust-p "I topi non avevano nipoti")
; -> T
\end{lstlisting}


\small\color{black!55}\hyperref[sol:139]{$\rightarrow$ Soluzione}\enspace\textcolor{black!30}{|}\enspace\hyperref[hint:139]{$\rightarrow$ Suggerimento}\normalsize\color{black}
\end{colbox}

\label{ex:140}
\begin{stdbox}{140}{Anagrammi}{3}
\exlabel{Implementare} \fn{(anagram-p s1 s2)}, ignorando maiuscole.
\nopagebreak
\begin{lstlisting}
(anagram-p "roma" "amor")   ; -> T
\end{lstlisting}


\small\color{black!55}\hyperref[sol:140]{$\rightarrow$ Soluzione}\enspace\textcolor{black!30}{|}\enspace\hyperref[hint:140]{$\rightarrow$ Suggerimento}\normalsize\color{black}
\end{stdbox}

\label{ex:141}
\begin{colbox}{141}{Conta caratteri}{2}
\exlabel{Implementare} \fn{(count-char c s)}.
\nopagebreak
\begin{lstlisting}
(count-char #\a "banana")   ; -> 3
\end{lstlisting}


\small\color{black!55}\hyperref[sol:141]{$\rightarrow$ Soluzione}\normalsize\color{black}
\end{colbox}

% =============================================================
\setcounter{section}{8}
\section{Parsing e interprete}
% =============================================================

\noindent
Un valutatore di espressioni simboliche senza usare \fn{eval}:
il codice come dato, pattern fondamentale di Lisp.

\bigskip

\label{ex:142}
\begin{stdbox}{142}{Calcolatrice RPN}{4}
\exlabel{Implementare} \fn{(rpn-eval espressione)}: valuta notazione
polacca inversa.
\begin{lstlisting}
(rpn-eval '(3 4 +))          ; -> 7
(rpn-eval '(3 4 * 2 +))      ; -> 14
(rpn-eval '(5 1 2 + 4 * + 3 -)) ; -> 14
\end{lstlisting}
Vincolo: usare una pila.



\small\color{black!55}\hyperref[sol:142]{$\rightarrow$ Soluzione}\enspace\textcolor{black!30}{|}\enspace\hyperref[hint:142]{$\rightarrow$ Suggerimento}\normalsize\color{black}
\end{stdbox}

\label{ex:143}
\begin{stdbox}{143}{Valutatore di espressioni}{4}
\exlabel{Implementare} \fn{(evaluate-expression expr)}: espressioni
in notazione prefissa, senza usare \fn{eval}.
\begin{lstlisting}
(evaluate-expression '(+ 2 3))        ; -> 5
(evaluate-expression '(* 4 (+ 2 3)))  ; -> 20
\end{lstlisting}


\small\color{black!55}\hyperref[sol:143]{$\rightarrow$ Soluzione}\enspace\textcolor{black!30}{|}\enspace\hyperref[hint:143]{$\rightarrow$ Suggerimento}\normalsize\color{black}
\end{stdbox}

\label{ex:144}
\begin{colbox}{144}{Mini-interprete con variabili}{5}
Estendere l'es.~143 con variabili (\fn{set x 10}) e riferimenti.
\begin{lstlisting}
(mio-eval '(set x 10) env)
(mio-eval '(+ x 5) env)   ; -> 15
\end{lstlisting}
Ambiente separato come alist; non usare \fn{eval}.



\small\color{black!55}\hyperref[sol:144]{$\rightarrow$ Soluzione}\enspace\textcolor{black!30}{|}\enspace\hyperref[hint:144]{$\rightarrow$ Suggerimento}\normalsize\color{black}
\end{colbox}

% =============================================================
\setcounter{section}{9}
\section{Progetti}
% =============================================================

\noindent
Progetti completi che integrano strutture dati, I/O, gestione
degli errori e organizzazione del codice.

\bigskip

\label{ex:145}
\begin{stdbox}{145}{Indovina il numero}{3}
Il programma sceglie un numero casuale tra 1 e 100,
l'utente lo indovina con feedback troppo alto/troppo basso.
\begin{lstlisting}
(gioca)
; Indovina il numero (1-100):
; > 50
; Troppo alto!
; ...
; Esatto! Ci hai messo 5 tentativi.
\end{lstlisting}
Vincolo: \fn{loop} per il ciclo; \fn{random} per la generazione.


\small\color{black!55}\hyperref[sol:145]{$\rightarrow$ Soluzione}\normalsize\color{black}
\end{stdbox}

\label{ex:146}
\begin{colbox}{146}{Tris (tic-tac-toe)}{4}
Gioco del tris a due giocatori da tastiera.
\begin{itemize}[leftmargin=2em]
  \item Griglia $3 \times 3$ con array.
  \item Verifica vittoria: righe, colonne, diagonali.
  \item Turno alternato; rilevamento pareggio.
\end{itemize}
Separare logica del gioco, verifica vittoria e I/O.



\small\color{black!55}\hyperref[sol:146]{$\rightarrow$ Soluzione}\enspace\textcolor{black!30}{|}\enspace\hyperref[hint:146]{$\rightarrow$ Suggerimento}\normalsize\color{black}
\end{colbox}

\label{ex:147}
\begin{colbox}{147}{Gioco della vita (Conway)}{4}
\exlabel{Implementare} \fn{(evolvi griglia)}: un passo di evoluzione.
Griglia: array 2D di 0 (morto) e 1 (vivo).
Regole: sopravvive con 2--3 vicini vivi; nasce con esattamente 3.
\begin{lstlisting}
(loop repeat 5
      do (setf g (evolvi g))
         (stampa-griglia g))
\end{lstlisting}


\small\color{black!55}\hyperref[sol:147]{$\rightarrow$ Soluzione}\enspace\textcolor{black!30}{|}\enspace\hyperref[hint:147]{$\rightarrow$ Suggerimento}\normalsize\color{black}
\end{colbox}

\label{ex:148}
\begin{stdbox}{148}{Mini sistema gestionale}{5}
Programma completo (studenti, libri, prodotti, film o contatti).

Requisiti minimi:
\begin{itemize}[leftmargin=2em]
  \item \fn{defstruct} o CLOS;
  \item hash table o lista come archivio;
  \item inserimento, ricerca, modifica, eliminazione,
        filtraggio, ordinamento;
  \item salvataggio e caricamento da file;
  \item gestione degli errori;
  \item docstring su ogni funzione pubblica.
\end{itemize}

Struttura minima:
\begin{lstlisting}
;; MODEL  -- strutture dati
;; LOGIC  -- operazioni
;; IO     -- input/output e file
\end{lstlisting}



\small\color{black!55}\hyperref[sol:148]{$\rightarrow$ Soluzione}\enspace\textcolor{black!30}{|}\enspace\hyperref[hint:148]{$\rightarrow$ Suggerimento}\normalsize\color{black}
\end{stdbox}







